\documentclass[11pt]{article}
\usepackage[a4paper,margin=2.2cm]{geometry}
\usepackage{amsmath,amssymb,amsthm}
\usepackage{array}
\usepackage{xcolor}
\usepackage{enumitem}
\usepackage[expansion=false]{microtype}
\usepackage[colorlinks=true,linkcolor=blue,citecolor=blue]{hyperref}

% ---------- symbols ----------
\newcommand{\rs}[1]{\lvert #1\rvert}                 % regex size
\newcommand{\rss}[1]{\lVert #1\rVert}                % size of a DISTINCT set (sum of sizes)
\newcommand{\Strong}{\mathsf{S}}                     % strong normaliser  (collapses a* a* -> a*)
\newcommand{\sig}{\sigma}                            % weak sequence plug (NO collapse)
\newcommand{\dpar}{\partial}                         % frontier (top alternation split)
\newcommand{\opn}{\mathrm{o}}                        % opened size of ONE regex
\newcommand{\front}{\mathrm{f}}                      % opened size of the frontier
\newcommand{\Dc}{\mathsf{D}}                         % Antimirov partial-derivative list
\newcommand{\Dch}{\widehat{\mathsf{D}}}              % canonicalised derivative rows
\newcommand{\Front}{\mathsf{F}}                      % afactored1: normalised front
\newcommand{\dstep}{\Delta}                          % rpder_strong_rows_raw: one strong step
\newcommand{\opens}{\Theta}                          % row_dlformss: open a row list
\newcommand{\dlf}{\mathrm{dl}}                       % row_dlforms: open one regex to linear forms
\newcommand{\sigt}{\sig_{\!7}}                       % rsimp7_SEQ_atom: sigma + top star-collapse
\newcommand{\Prune}{\mathsf{P}}                      % rsimpStrong_prune_rows_raw
\newcommand{\pot}{\mathsf{p}}                        % strong_opened_live_acc_potential
\newcommand{\SOL}{\mathsf{U}}                        % strong_opened_live_row_universe
\newcommand{\tr}{\tau}                               % rstar_atrace: the event list
\newcommand{\wt}{\mathsf{w}}                          % aevt_cost: event cost
\newcommand{\seq}{\cdot}
\newcommand{\alt}{+}
\newcommand{\one}{\mathbf 1}
\newcommand{\zero}{\mathbf 0}
\newcommand{\eps}{\varepsilon}
\newcommand{\GREEN}{\textcolor{green!55!black}{\textbf{[proved]}}}
\newcommand{\VALID}{\textcolor{orange!85!black}{\textbf{[validated, not yet formalised]}}}
\newcommand{\OPEN}{\textcolor{red}{\textbf{[OPEN]}}}
\newcommand{\WHY}{\textit{\textcolor{blue}{Why:\ }}}

\theoremstyle{definition}
\newtheorem*{goal}{Goal}
\setlength{\parindent}{0pt}\setlength{\parskip}{4pt}

\title{\textbf{The cubic size bound: current proof state}\\[2pt]
\large the formalised chain, the open envelope and its validated scalar crux, and why it resists}
\author{Secretary working note --- all numeric facts machine-validated}
\date{2026-06-16}

\begin{document}\maketitle

\section{Symbols (every function defined once, recursively)}\label{sec:notation}

Regexes: $\zero,\one,$ char $c$, sequence $r_1\seq r_2$, alternation $\sum_i r_i$, star $r^{*}$. An \emph{atom} $a$ is a
char, alternation or star (non-$\seq$, non-unit).

\begin{description}[leftmargin=2.2em,style=nextline,itemsep=3pt]

\item[$\rs r$ \rm(size)] $\rs\zero=\rs\one=\rs c=1,\ \rs{r_1r_2}=1+\rs{r_1}+\rs{r_2},\ \rs{\sum rs}=1+\sum\rs{r_i},\ \rs{r^{*}}=1+\rs r$.

\item[$\rss X$ \rm(set size)] $\sum_{x\in X}\rs x$ over the \emph{distinct} members of a finite set/list $X$.

\item[$\Strong$ \rm(strong normaliser; ``rsimpStrong\_raw'')] simplify bottom-up, plugging sequences with $\sigt$ (so the
guarded $a^{*}\seq a^{*}\to a^{*}$ fires at every level), and idempotent-collapse stars:
\[
\Strong\zero=\zero,\ \ \Strong\one=\one,\ \ \Strong c=c,\qquad
\Strong(r_1r_2)=\sigt(\Strong r_1,\,\Strong r_2),\qquad
\Strong(\textstyle\sum rs)=\mathsf A\bigl(\mathrm{flat}(\mathrm{map}\,\Strong\,rs)\bigr),
\]
\[
\Strong(r^{*})=\bigl(\text{case }\Strong r\text{ of }\ \zero\Rightarrow\one\ \mid\ \one\Rightarrow\one\ \mid\ s^{*}\Rightarrow s^{*}\ \mid\ s\Rightarrow s^{*}\bigr)
\]
($\mathsf A$ flattens, dedupes and rebuilds an alternation). The single collapse $a^{*}\seq a^{*}\to a^{*}$ --- realised by
$\sigt$ in the sequence case, recursively --- is the source of all the difficulty below.

\item[$\sig(h,k)$ \rm(weak plug: ``$h$ then $k$'')] does all sequencing; \textbf{never} collapses $a^{*}\seq a^{*}$:
\begin{gather*}
\sig(\zero,k)=\zero,\quad \sig(\one,k)=k,\quad \sig(r_1\seq r_2,k)=\sig\!\bigl(r_1,\sig(r_2,k)\bigr),\\
\sig(a,\zero)=\zero,\quad \sig(a,\one)=a,\quad \sig(a,k)=a\seq k.
\end{gather*}
$\sig(p,\one)$ \emph{canonicalises} $p$ alone (right-associate, erase units, absorb $\zero$; e.g.\ $\sig((ab)c,\one)=a(bc)$,
$\sig(a\one,\one)=a$).

\item[$\sigt(h,k)$ \rm(strong plug: ``rsimp7\_SEQ\_atom'')] $\sig$ \emph{plus} the guarded $a^{*}\seq a^{*}\to a^{*}$
collapse applied \textbf{only at the top head}: $\sigt(a^{*},a^{*})=a^{*}$, $\sigt(a^{*},a^{*}\seq k)=a^{*}\seq k$, and
$\sigt(h,k)=\sig(h,k)$ otherwise. So $\sigt$ sits between $\sig$ (no collapse) and $\Strong$ (collapses everywhere) ---
it collapses just the leading star pair.

\item[$\dpar k$ \rm(frontier)] split the top alternation of $k$ into its branches ($\dpar\zero=\emptyset$,
$\dpar(\sum rs)=\bigcup\dpar r_i$, else $\{k\}$).

\item[$\dlf(x)$ \rm(open one regex to linear forms; ``row\_dlforms'')] by cases on the top constructor:
\begin{gather*}
\dlf(\zero)=\emptyset,\quad \dlf(\textstyle\sum rs)=\textstyle\bigcup_{q}\dlf(q),\quad \dlf(r)=\dpar r\ \ (\text{otherwise}),\\
\dlf\bigl((\textstyle\sum ps)\seq k\bigr)=\textstyle\bigcup_{p\in ps}\dlf(\sigt(p,k)).
\end{gather*}
Distribute a top alternation, and an alternation-headed sequence's branches over the tail via $\sigt$; bottom out at the
frontier. (This is the \emph{only} other place $a^{*}\seq a^{*}$ can collapse --- via $\sigt$ in the $\seq$ case.)

\item[$\opn(x)$, $\front(k)$ \rm(opened size of one regex / of a frontier)]
$\opn(x):=\rss{\dlf(\Strong x)}$ (strong-simplify, open, dedupe, sum sizes);\quad
$\front(k):=\rss{\bigcup_{q\in\dpar k}\dlf(\Strong q)}$.

\item[$\Dc_c(r)$ \rm(Antimirov derivative list)] the partial-derivative ``rows'', sequencing by $\sig$:
\[
\Dc_c\zero=\Dc_c\one=[],\quad \Dc_c(\text{char }d)=[\one\text{ if }d{=}c\text{ else }],\quad \Dc_c(\textstyle\sum rs)=\textstyle\bigoplus_i\Dc_c(r_i),
\]
\[
\Dc_c(r_1r_2)=[\sig(p,r_2)\mid p\in\Dc_c(r_1)]\ \oplus\ (\Dc_c(r_2)\text{ if }r_1\text{ nullable}),\qquad
\Dc_c(r^{*})=[\sig(p,r^{*})\mid p\in\Dc_c(r)].
\]

\item[$\Dch_c(r)$ \rm(canonicalised rows)] $:=[\sig(p,\one)\mid p\in\Dc_c(r)]$ --- every row put in canonical form so the
dedup below actually merges associativity/unit variants (keeping the row set small).

\item[$\Front(r,s)$ \rm(normalised front; ``afactored1'')] read the whole string $s$ with \emph{no} $\Strong$, threading the
\emph{row-list} (never collapsing it between letters): starting from $[r]$, apply for each letter $c$ the normalised step
$\nu_c(\rho):=\mathrm{dedup}\bigl(\mathrm{flat}\bigcup_{q\in\rho}\Dch_c(q)\bigr)$. Writing $G$ for the list-stepping (``afactored\_steps''),
\[
G(\rho,\eps)=\rho,\qquad G(\rho,\,c\!:\!s)=G\bigl(\nu_c(\rho),\,s\bigr),\qquad
\Front(r,s):=G([r],s)\ =\ \nu_{c_n}\!\circ\cdots\circ\nu_{c_1}([r])\quad(s=c_1\cdots c_n).
\]

\item[$\dstep_c(\rho)$ \rm(one strong step; ``rpder\_strong\_rows\_raw'')] derive+canonicalise each row, strong-simplify with
$\Strong$, flatten, strong-prune with $\Prune$, flatten, dedupe:
\[
\dstep_c(\rho):=\mathrm{dedup}\Bigl(\mathrm{flat}\;\Prune\bigl(\mathrm{flat}\textstyle\bigcup_{q\in\rho}[\,\Strong(x)\mid x\in\Dch_c(q)\,]\bigr)\Bigr).
\]

\item[$\Prune(\rho)$ \rm(strong prune; ``rsimpStrong\_prune\_rows\_raw'')] $\Prune(\rho):=\Prune_{\!*}([],\rho)$ --- a
left-to-right scan with an accumulator $\mathit{sn}$ of the rows already pruned:
\[
\Prune_{\!*}(\mathit{sn},[])=[],\qquad
\Prune_{\!*}(\mathit{sn},\,r\!:\!\rho)=\bigl(\text{let }r'=\pi^{*}(\mathit{sn},r)\ \text{in}\ r'\#\Prune_{\!*}(r'\!:\!\mathit{sn},\rho)\bigr).
\]
Prune one row against all earlier ones, then the pairwise rule (sharing a right factor $k$):
\[
\pi^{*}([],r)=r,\qquad \pi^{*}(x\!:\!xs,\,r)=\pi^{*}\bigl(xs,\,\pi(x,r)\bigr),
\]
\[
\pi(e,l)=\begin{cases}\sigt\bigl(\textstyle\sum(rrs\setminus lrs),\,k\bigr)&\text{if }\ e=(\textstyle\sum lrs)\seq k,\ \ l=(\textstyle\sum rrs)\seq k,\\[2pt] l&\text{otherwise.}\end{cases}
\]
where $rrs\setminus lrs$ keeps the branches of $rrs$ not occurring in $lrs$ ($[]\mapsto[]$; $r\!:\!rs\mapsto$ drop $r$ if
$r\in lrs$, else keep $r$), and $\textstyle\sum$ rebuilds the alternation. So a row is \emph{shrunk, never dropped}
(length-preserving); the rule is language-preserving and size-non-increasing.

\item[$\opens(\rho)$ \rm(open a row list; ``row\_dlformss'')] $\opens(\rho):=\bigcup_{q\in\rho}\dlf(q)$.

\item[$\pot(r,k)$ \rm(potential)] an over-approximation of the opened size of $r$'s derivative universe against $k$,
defined to recurse on $r$ and bound it from above ($\rss{\SOL(r)}\le\pot(r,\one)$, proved):
\[
\pot(\zero,k)=1,\quad \pot(\one,k)=1+\opn(k)+\front(k),\quad \pot(\text{char }c,k)=\opn(\sig(c,k))+1+\opn(k)+\front(k),
\]
\[
\pot(r_1r_2,k)=\pot(r_1,\sig(r_2,k))+\pot(r_2,k),\quad
\pot(\textstyle\sum rs,k)=1+\opn(\sig(\textstyle\sum rs,k))+\textstyle\sum_{q}\pot(q,k),
\]
\[
\pot(r^{*},k)=\underbrace{\opn(\sig(r^{*},k))}_{\text{head}}+\underbrace{\pot\bigl(r,\sig(r^{*},k)\bigr)}_{\text{body: $r$ against its own star}} .
\]

\item[$\SOL(r)$ \rm(strong-opened universe; ``strong\_opened\_live\_row\_universe'')] open (via $\dlf\circ\Strong$) every
regex of $r$'s derivative universe $\mathcal D(r)$:
\[
\SOL(r):=\textstyle\bigcup_{p\in\mathcal D(r)}\dlf(\Strong p),\qquad
\mathcal D(r):=\{\zero,\one,r\}\,\cup\,\mathrm{paths}(r)\,\cup\,\dpar r\,\cup\,\textstyle\bigcup_{q\in\mathrm{paths}(r)}\dpar q,
\]
where $\mathrm{paths}(r)$ is the (finite) set of continuations reachable by repeatedly deriving $r$. Proved static facts:
$\#\mathcal D(r)\le\rs r+2$ (linear) and each member has size $\le(\rs r+2)^2+1$ (quadratic) --- the linear$\times$quadratic
that the cubic budget rests on. $\rss{\SOL(r)}$ is what the Gate must bound.

\item[$\tr(r)$ \rm(trace), $\wt_k$ \rm(event cost)] the explicit event-list expansion of $\pot(r^{*},\cdot)$, with the
continuation tracked as a prefix stack $xs$ that $\mathrm{ap}$ instantiates:
\[
\mathrm{ap}(a,[],k)=\sig(a^{*},k),\qquad \mathrm{ap}(a,\,x\!:\!xs,\,k)=\sig\bigl(x,\mathrm{ap}(a,xs,k)\bigr).
\]
Events are $\mathrm{AU}$ (unit), $\mathrm{AO}\,xs$ (open), $\mathrm{AF}\,xs$ (frontier), with cost
$\wt_k(\mathrm{AU})=1$, $\wt_k(\mathrm{AO}\,xs)=\opn(\mathrm{ap}(r,xs,k))$,
$\wt_k(\mathrm{AF}\,xs)=\rss{\bigcup_{q\in\dpar(\mathrm{ap}(r,xs,k))}\dlf(\Strong q)}$. The trace $\mathrm{at}$ mirrors
$\pot$, pushing the right part onto the stack:
\begin{gather*}
\mathrm{at}(\zero,xs)=[\mathrm{AU}],\qquad \mathrm{at}(\one,xs)=[\mathrm{AU},\mathrm{AO}\,xs,\mathrm{AF}\,xs],\\
\mathrm{at}(c,xs)=\mathrm{AO}(c\!:\!xs)\#[\mathrm{AU},\mathrm{AO}\,xs,\mathrm{AF}\,xs],\qquad
\mathrm{at}(pq,xs)=\mathrm{at}(p,\,q\!:\!xs)\oplus\mathrm{at}(q,xs),\\
\mathrm{at}(\textstyle\sum rs,xs)=\mathrm{AU}\#\mathrm{AO}(\textstyle\sum rs\!:\!xs)\#\textstyle\bigoplus_q\mathrm{at}(q,xs),\\
\mathrm{at}(r^{*},xs)=\mathrm{AO}(r^{*}\!:\!xs)\#\mathrm{at}(r,\,r^{*}\!:\!xs);
\end{gather*}
$\tr(r):=\mathrm{AO}\,[]\ \#\ \mathrm{at}(r,[])$, with the \emph{proved identity} $\pot(r^{*},k)=\sum_{e\in\tr(r)}\wt_k(e)$.

\item[\rm minor list operations] $\mathrm{flat}$ flattens nested alternations and drops $\zero$s; $\mathrm{dedup}$ keeps
the first occurrence of each row; $\oplus$/$\bigoplus$ is list append/concat; ``$r$ nullable'' means $r$ matches
$\eps$; $rrs\setminus lrs$ drops from $rrs$ the entries already occurring in $lrs$.

\end{description}

\section{The goal, and why it is cubic}

In the clean fragment (classical regexes, normal form, no bounded repetition): read $s$, take \emph{one} strong step on
the next letter $c$, open the resulting rows and sum their sizes. The thesis target is that this is cubic in $\rs r$,
uniformly in $s$:
\[
\boxed{\ \rss{\,\opens\!\bigl(\dstep_c(\Front(r,s))\bigr)\,}\ \le\ 2\,(\rs r+3)^3\ }
\tag{Gate}
\]
\WHY a cubic ceiling independent of $\rs s$ makes the lexer's space polynomial in the regex regardless of input length;
it is the crux. \emph{Design note:} $\Front(r,s)$ is the cheap $\Strong$-free accumulation (the D~law and the static cubic
front are proved on it); the $\Strong$ collapse $a^{*}\seq a^{*}\to a^{*}$ and the pruning are applied \emph{once}, only in
the final $\dstep_c$. So $\Strong$'s collapse --- and the global cancellation it induces --- live only at that last step,
and that is the wall (\S\ref{sec:open}--\S\ref{sec:why}).

\section{The chain (top to bottom)}\label{sec:chain}

Each line is justified by the one below it. Steps 1--6 and 9 are machine-checked (no \texttt{sorry}); the whole chain is
formalised \emph{modulo one hypothesis}, the affine envelope $\mathrm{env}$ of step~4, carried explicitly by the gate lemma
\texttt{actual\_gate\_from\_cube\_shell}. Steps 7--8 --- the reduction of $\mathrm{env}$ to the single scalar ($\star$) ---
are machine-validated but \textbf{not yet formalised} (see the colour key).

\begin{enumerate}[leftmargin=1.6em,itemsep=5pt]

\item \GREEN\ \textbf{Gate $\Leftarrow$ root potential.} The next-step opened rows lie inside $\SOL(r)$ (monotonicity),
and $\rss{\SOL(r)}\le\pot(r,\one)$ is proved, so the Gate follows from $\pot(r,\one)\le 2(\rs r+3)^3$.
\WHY trade a statement about an opened \emph{set} for a recursively-defined numeric \emph{upper bound} $\pot$ that
recurses cleanly on the regex.

\item \GREEN\ \textbf{The cube-shell.} For all $r,k$:\ $\pot(r,k)\le(\rs r+\rs k)^3-(\rs k)^3$ (CS); at $k=\one$,
$\pot(r,\one)\le(\rs r+1)^3-1\le2(\rs r+3)^3$, closing step 1.
\WHY (CS) is a difference of cubes --- it telescopes additively, hence is provable by induction on $r$.

\item \GREEN\ (except one case) \textbf{The cube-shell induction (on $r$).} $\zero,\one,$ char are leaves;
$\sum$ and $r_1r_2$ compose the children's shells --- \emph{all proved}. The one open constructor is $r^{*}$.
\WHY the $r^{*}$ recursion feeds the body $r$ against an \emph{inflated} continuation (step 5), so the naive use of (CS)
as the body's hypothesis overshoots by ${\approx}7\times$ and cannot close that case.

\item \GREEN\ \textbf{$r^{*}$ case $\Leftarrow$ the affine envelope.} With $M:=\rs{r^{*}}=1+\rs r$, it suffices to prove
\[
\pot(r^{*},k)\ \le\ M^3+3M^2\,\rs k.\tag{Env}
\]
Indeed $M^3+3M^2\rs k\le(M+\rs k)^3-(\rs k)^3$ (leftover $3M(\rs k)^2\ge0$), so (Env)$\Rightarrow$(CS) at $r^{*}$.
\WHY (Env) is \emph{linear in $\rs k$}, where (CS) is quadratic; (Env) is the tight statement at $r^{*}$. It is
machine-validated true ($0$ violations).

\item \GREEN\ \textbf{Unfolding $r^{*}$.} $\pot(r^{*},k)=\opn(\sig(r^{*},k))+\pot(r,\sig(r^{*},k))$. The head is harmless;
the \textbf{body} opens $r$ against $\sig(r^{*},k)=r^{*}\seq k$, its own star re-entry, with continuation inflated by ${\approx}M$.
\WHY the heart of the difficulty: the body is opened against a continuation containing a copy of its own star. The opened
size does \emph{not} grow with the inflation (the re-entry saturates via $a^{*}\seq a^{*}\to a^{*}$, used helpfully), but
capturing that is the open problem.

\item \GREEN\ \textbf{(Env) $\Leftarrow$ a per-event certificate.} By the trace identity $\pot(r^{*},k)=\sum_{e\in\tr(r)}\wt_k(e)$.
Given $a_1,a_B:\text{events}\to\mathbb N$ with
\begin{gather*}
\text{(affine)}\ \wt_k(e)\le a_1(e)+a_B(e)(\rs k-1),\\
\text{(slope)}\ \textstyle\sum_e a_B(e)\le 3M^2,\qquad
\text{(intercept)}\ \textstyle\sum_e a_1(e)\le M^3+3M^2,
\end{gather*}
the certificate proves (Env) --- lemma \texttt{RSTAR\_affine\_envelope\_from\_certificate}, proved but currently
\emph{unused}: it is never instantiated, because its slope/affine hypotheses are exactly steps 7--8, which are not yet
formalised. So the gate still carries $\mathrm{env}$ as an open hypothesis.
\WHY splits (Env) into a continuation-slope and a continuation-free intercept, each a single sum.

\item \VALID\ \textbf{Slope.} $a_B$ a row/frontier count per event should give $\sum_e a_B(e)\le3M^2$.
\WHY $k$ sits at the \emph{tail} of every opened row (behind the single anchor $r^{*}$), so it adds size linearly.
\emph{Validated, but no concrete $a_B$ and no slope lemma exist in the source yet.}

\item \VALID\ \textbf{Intercept reduced to a single anchor.} Take $a_1(e):=\wt_{\,\text{char }c}(e)$ for a fixed char $c$:
one fixed tail dominates every size-$1$ tail per event (so (affine) holds), and then
$\sum_e a_1(e)=\pot(r^{*},\text{char }c)$ by the trace identity --- collapsing the intercept to ($\star$).
\WHY dissolves a multiplicity trap (a per-event max over all size-$1$ tails over-counts); one fixed tail collapses the
intercept sum back to one potential value. \emph{Validated; the per-event affine domination is not yet a formal lemma.}

\item \GREEN\ \textbf{Event count is linear.} $\rs{\tr(q)}\le 4\rs q$ (length of the trace).
\WHY the linear factor: the intercept is \emph{(linearly many events)} $\times$ \emph{(per-event opened size)}, and we
must show this product is cubic.

\end{enumerate}

\section{The open piece: the envelope, and its validated scalar crux}\label{sec:open}

\textbf{What the formalised chain actually carries open} is the affine envelope $\mathrm{env}$ (step 4), a single hypothesis
of the gate lemma, quantified over \emph{every} clean star sub-term and \emph{every} continuation:
\[
\mathrm{env}:\quad \pot(r^{*},k)\ \le\ M^3+3M^2\,\rs k\qquad(\text{all clean }r^{*},\text{ all }k),\quad M=1+\rs r.
\]
The proved-but-unused certificate (step 6) shows $\mathrm{env}\Leftarrow(\text{slope})\wedge(\text{affine})\wedge(\text{intercept})$.
Steps 7--8 (\VALID) supply the slope and pick the single fixed tail, collapsing the intercept to the \emph{anchor} ---
$\mathrm{env}$ at its hardest instance $\rs k=1$:
\[
\boxed{\ \pot(r^{*},\,\text{char }c)\ \le\ M^3+3M^2\ }
\qquad\Bigl(=\ \textstyle\sum_{e\in\tr(r)}\wt_{\,\text{char }c}(e)\ \le\ M^3+3M^2\Bigr),\quad M=1+\rs r.
\tag{$\star$}
\]
\OPEN\ (scalar), and the bridge $(\star)+\text{slope}+\text{affine}\Rightarrow\mathrm{env}$ still needs formalising (steps 7--8).
($\star$) is machine-validated true: $0$ violations over an \emph{exhaustive} enumeration of all in-regime terms up to
$\rs r=9$ ($\approx8.7\times10^5$ cases); worst ratio $0.95$ at $M=2$ ($19\le20$), slacker for larger $M$ (true value
${\approx}0.5\times$ budget). Unfolding $r^{*}$, ($\star$) is the \emph{body saturation}
$\pot(r,\sig(r^{*},\text{char }c))\le M^3+3M^2-\opn(\sig(r^{*},\text{char }c))$: the body opened against its own star
re-entry, budgeted at the \emph{original} (size-$1$) tail, not the inflated one.

\section{Why ($\star$) resists --- a global cancellation}\label{sec:why}

($\star$) is $\sum_{\text{events}}(\text{opened size})\le M^3+3M^2$, with the event count linear (step 9) and each opened
size $\le$ quadratic. Naively ``linear $\times$ quadratic $=$ cubic'' --- yet every \emph{structural} bound (holding
pointwise at each event/node/row) overshoots, while the true sum is ${\approx}0.5\times$ the budget:

\begin{center}\small
\begin{tabular}{ll@{\quad}>{\raggedright\arraybackslash}p{5.6cm}}
\textbf{candidate bound} & \textbf{ratio to budget} & \textbf{fate}\\\hline
true value $\sum_e\wt_{\,c}(e)$ & ${\approx}0.5$ & what we must bound\\
crude (events $\times$ max event cost) & ${\approx}1.5\times$ & overshoots\\
saturated per-event cap $(M{+}1)M\times4\rs r$ & ${\approx}4\times$ & overshoots\\
generic quadratic opening $\mathrm{Suc}(sz)\cdot sz$ & ${\approx}7\times$ & overshoots\\
tight body invariant $(\rs q{+}1)(M{+}1)M$ & $1.0$ pointwise & does \emph{not} telescope ($\sum$: $1+\text{cap}\cdot|rs|\le\text{cap}$ false)\\
deduped universe $\rss{\text{rows}(r^{*})}$ & ${\approx}0.2$ (cubic!) & per-row multiplicity reaches $18$; no injection\\
\end{tabular}
\end{center}

\textbf{Diagnosis.} The sum is small only by a \textbf{global cancellation}: events opening \emph{many} rows have
\emph{small} per-row cost; events with \emph{few} rows have \emph{large} cost; high-multiplicity rows are small. Size and
multiplicity are \emph{anti-correlated}. A structural induction must bound the worst case \emph{at every event at once},
so it charges high multiplicity and high size simultaneously and overshoots. \textbf{No \texttt{rsize}/count
structural-induction invariant sees this cancellation.} (Hence the clean ``deduped universe is cubic'', ratio $0.2$,
does not help: the trace charges some rows $18\times$, and the multiplicity-weighted sum has no injection into the
deduped set.) On exactly ($\star$): $8{+}$ external rounds, a $4$-agent design panel (a confident \emph{false} bound,
caught), and three formalisation attempts --- all structural, all overshoot.

\section{The shape of a correct proof: the amortised precedent}\label{sec:amort}

The sibling problem --- the row \emph{count} being linear (the ``D~law'') --- had the same flavour and \emph{was} cracked,
not by structural induction but by a \textbf{simultaneous induction on two mutually-supporting invariants}: a
boundary-strengthened bound $T$ and a strict credit $S$ (with $\mathrm{acc}$ the accumulator, $\mathrm{zw}(r)=\#\text{letters}+\#\text{stars}\le\rs r$):
\[
T:\ \#\bigl((\dpar\sig(r,k)\cup\mathrm{acc}(r,k))\setminus\dpar k\bigr)\le\mathrm{zw}(r),
\quad
S:\ \mathrm{zw}(r)>0\Rightarrow 1+\#\bigl(\mathrm{acc}(r,k)\setminus\dpar k\bigr)\le\mathrm{zw}(r).
\]
$S$ is the amortisation: a constructor may add a boundary row no \emph{local} cap can absorb, but it is paid by the
strict spare credit carried in the strengthened $T$, and $T\wedge S$ (one induction) gives the count bound.
\WHY a global ``adds-but-pays'' cancellation is captured by carrying a credit in a strengthened invariant, not by a
pointwise bound.

\textbf{The analogue we need (open).} A \emph{size} version of $T+S$: a potential $\Phi$ on the state
(root~$r^{*}$, sub-term~$q$, boundary) plus a strict \emph{size}-credit, so the per-event opened size
is paid by the drop of $\Phi$ across the trace, telescoping to $\le M^3+3M^2$ at the root, and surviving (i) the $\sum$
case (where pointwise caps fail) by spending the credit once, and (ii) the star re-entry by the boundary being idempotent
at the anchor. No concrete such $\Phi$ has yet been exhibited.

\section{Where the proof can go}

\begin{enumerate}[leftmargin=1.6em,itemsep=3pt]
\item \textbf{Build the amortised $\Phi$ (\S\ref{sec:amort}) concretely} and verify it telescopes and stays $\le M^3+3M^2$
--- the direct analogue of the (successful) count proof; the most promising untried route.
\item \textbf{Rethink the opening/$\pot$ definitions} so the cancellation becomes a \emph{local} invariant (then a
structural induction suffices) --- it may be the honest answer if no $\Phi$ on the current $\pot$ is local.
\item \textbf{Treat ($\star$) as a validated lemma} --- exhaustively checked ($0/8.7\times10^5$); the rest of the Gate is
then proved \emph{modulo formalising} the (also validated) slope and single-tail affine steps that lift ($\star$) to the
full envelope $\mathrm{env}$ --- and the certificate that consumes them is already proved.
\end{enumerate}

\paragraph{One-line status.} The clean-fragment cubic Gate is formalised \emph{modulo the affine envelope} $\mathrm{env}$;
a proved (but as-yet unused) certificate reduces $\mathrm{env}$ --- via a machine-validated slope bound and single-tail
affine step that still need formalising --- to the single scalar $\pot(r^{*},\text{char }c)\le M^3+3M^2$, itself
exhaustively validated and provably beyond structural induction (a global size$\times$multiplicity cancellation). The route
that fits is the size analogue of the $T+S$ amortised argument that cracked the row count.

\end{document}
